\chapter{Choosing}\label{ch:in:choosing}

Four offers, a bank quant role, a market maker's trading role, an analyst seat at a platform and a technology firm's
engineering job, differ in expected pay, in its spread, in what tax and housing take, in hours, in what the next five
years teach and in how likely the job is to last. No single number ranks them: a person's weights do, and weights are
hard to state. The useful questions are which offers no weights can save, how often each offer comes first when the
weights are uncertain, and how far a weight must move before the choice changes. This chapter builds that analysis
from the tools of the previous chapters and closes the book.

\section{The attributes that differ across firm types}

The chapters of Parts~II and~III measured what differs.
\begin{itemize}
\item \textbf{Pay}, its structure and its spread (chapters~13 and~14): base and bonus, deferral and \omterm{def:in:how-pay-works:rsu}{forfeiture},
formulaic payouts, and the ranges the filings show.
\item \textbf{Place} (chapters~15 and~27): tax, housing and price levels.
\item \textbf{Hours} (chapter~26): the coverage a desk needs and the on-call it shares.
\item \textbf{The work} (chapters~16 to~25): what a role ships, how fast it learns whether it is good (chapter~17's
\omterm{def:in:quant-researcher:feedback}{feedback speed}) and what it leads to (chapter~28).
\item \textbf{Security}: the stops of a platform (chapter~22) and the chance that a job ends for reasons the holder does
not control.
\end{itemize}

\begin{definition}[Multi-attribute value model, swing weighting]\label{def:in:choosing:mavm}
A \emph{multi-attribute value model}\index{multi-attribute value model} scores each alternative on each attribute, scales
every attribute from the worst to the best alternative, and adds the scaled scores with weights (Keeney and Raiffa).
\emph{Swing weighting}\index{swing weighting} sets the weights by asking how much the chooser values moving each
attribute from its worst to its best level, relative to the others, rather than how important the attribute is in the
abstract.
\end{definition}

\section{Pay risk as a certainty equivalent}

Chapter~13 turned a package's distribution into a certainty equivalent: the sure amount with the same expected utility
under constant relative risk aversion (Book~4, chapter~9). With relative risk aversion 3 and \$500\,000 of other wealth,
the chapter's four offers are worth \$1\,186\,354 (bank quant), \$1\,365\,186 (market maker), \$903\,603 (platform
analyst) and \$845\,964 (technology engineer) over five years. The certainty equivalent folds risk into pay; the value
model below keeps them apart, as expected pay and pay risk (the 10th--90th percentile range over the mean), so that the
weight on risk can be seen and varied.

\begin{center}\footnotesize\setlength{\tabcolsep}{4pt}
\begin{tabular}{@{}lrrrrrr@{}}
\toprule
offer (city) & expected pay & pay risk & tax, housing & on-call & learning & security\\
 & \$k, 5 years & range/mean & share & nights/month & 1--5 & share\\
\midrule
bank quant (London) & 1\,618 & 1.01 & 53.1\% & 0.0 & 3 & 1.00\\
market maker (Chicago) & 1\,949 & 1.16 & 40.8\% & 3.8 & 4 & 1.00\\
platform analyst (New York) & 1\,703 & 1.67 & 50.0\% & 0.0 & 4 & 0.65\\
tech engineer (London) & 1\,191 & 1.22 & 55.3\% & 5.1 & 3 & 1.00\\
\bottomrule
\end{tabular}
\end{center}

The table's pay figures come from chapter~13's offers and a fourth built the same way (a technology firm's base with a
restricted-stock grant vesting over four years); every package is illustrative. The share taken by tax and housing
applies chapter~27's tools to each offer's average annual pay in its city; on-call nights use chapter~26's rotas
(eight people for the market maker's round-the-clock desk, six for the engineer's service); learning is the chooser's
judgement; security is the share of five-year careers in which the pay model does not close the job, which only the
platform's book can do.

\section{Temperament, skills and the work itself}

Some attributes resist numbers and matter most. Chapter~17's \omterm{def:in:quant-researcher:feedback}{feedback speed} decides how soon a person learns whether
they are good at the job; chapter~16's supervision load and chapter~26's hours decide what a day is like; chapter~22's
stops decide how much a bad year costs. A person who dislikes uncertain pay should say so in the weight on pay risk; a
person who wants to build systems rather than trade them should say so in the learning score. The model does not know
temperament; it makes the chooser state it and shows what follows.

\section{A multi-attribute framework and its sensitivity}

\begin{definition}[Dominated alternative, rank acceptability]\label{def:in:choosing:smaa}
An alternative is a \emph{dominated alternative}\index{dominated alternative} when another is at least as good on every
attribute and better on one: no weights can make it the best. The \emph{rank acceptability}\index{rank acceptability}
of an alternative for a rank is the share of weight vectors, drawn from a stated distribution, under which the
alternative takes that rank (Lahdelma, Hokkanen and Salminen's stochastic multicriteria acceptability analysis).
\end{definition}

\begin{method}[Screening, weighting and sensitivity]\label{met:in:choosing:method}
Scale each attribute to $[0,1]$ from the worst to the best offer (reversed where less is better). Drop dominated offers.
With swing weights $w$, the value of offer $i$ is $\sum_j w_j s_{ij}$. Draw weights uniformly on the simplex (a
Dirichlet distribution with unit parameters) and count how often each offer takes each rank. For the attribute of
interest, find the smallest change of its weight, the others rescaled in proportion, that changes the best offer.
\end{method}

With the illustrative swing points (expected pay 30, pay risk 15, place 10, on-call 15, learning 20, security 10), the
market maker's offer scores 0.852, the platform's 0.589, the bank's 0.584 and the technology engineer's 0.202. The
technology offer is dominated: the bank's offer pays more with less spread, takes less in tax and housing, asks for no
on-call, and matches it on learning and security. Under uniformly random weights (\cref{fig:in:choosing:ranks}) the
market maker comes first in 79.5\% of draws, the bank in 15.9\%, the platform in 4.6\% and the technology offer never.

\begin{omfigure}
\begin{tikzpicture}
\begin{axis}[xbar stacked,width=10.5cm,height=5cm,xmin=0,xmax=100,ymin=0.4,ymax=4.6,bar width=12pt,ytick=data,
  yticklabels from table={figdata/industry/30-choosing/ranks.csv}{offer},yticklabel style={font=\footnotesize},
  y dir=reverse,xticklabel style={font=\footnotesize},xlabel={\small share of weight draws, \%},table/col sep=comma,
  legend style={legend cell align=left,font=\scriptsize,at={(0.5,-0.3)},anchor=north,legend columns=4,/tikz/every even column/.append style={column sep=8pt}}]
\addplot[fill=omDef,draw=none] table[x=r1,y=pos]{figdata/industry/30-choosing/ranks.csv};
\addplot[fill=omDef!55,draw=none] table[x=r2,y=pos]{figdata/industry/30-choosing/ranks.csv};
\addplot[fill=omThm!45,draw=none] table[x=r3,y=pos]{figdata/industry/30-choosing/ranks.csv};
\addplot[fill=black!25,draw=none] table[x=r4,y=pos]{figdata/industry/30-choosing/ranks.csv};
\legend{first,second,third,fourth}
\end{axis}
\end{tikzpicture}

\omcaption{\omterm{def:in:choosing:smaa}{Rank acceptability} of the four illustrative offers: the share of 20\,000 weight vectors, drawn uniformly on
the simplex of six attributes, under which each offer takes each rank. Data: \texttt{in\_choose.results}, from
\texttt{firm.careerdec.rank\_acceptability}.}\label{fig:in:choosing:ranks}
\end{omfigure}

The market maker's lead is robust to the weight on pay risk: that weight must rise by 0.452, from 0.15 to 0.602 with the
other weights shrunk in proportion, before the bank's offer becomes the best (\cref{fig:in:choosing:sweep}). A chooser
for whom the spread of pay outweighs expected pay, place, hours, learning and security together should take the bank's
offer; anyone else, on these numbers, the market maker's.

\begin{omfigure}
\begin{tikzpicture}
\begin{axis}[width=10.5cm,height=5.6cm,xmin=0,xmax=0.95,ymin=0,ymax=1,xlabel={\small weight on pay risk},
  ylabel={\small value},tick label style={font=\footnotesize},grid=major,grid style={gray!20},table/col sep=comma,
  legend style={font=\scriptsize,at={(0.5,-0.3)},anchor=north,legend columns=2,legend cell align=left,
  /tikz/every even column/.append style={column sep=8pt}}]
\addplot[omDef,thick] table[x=w,y=bank_quant]{figdata/industry/30-choosing/sweep.csv};
\addplot[omThm,thick] table[x=w,y=market_maker]{figdata/industry/30-choosing/sweep.csv};
\addplot[black!60,thick,dashed] table[x=w,y=platform_analyst]{figdata/industry/30-choosing/sweep.csv};
\addplot[black!30,thick] table[x=w,y=technology_engineer]{figdata/industry/30-choosing/sweep.csv};
\draw[black!50,dashed] (axis cs:0.15,0) -- (axis cs:0.15,1);
\draw[black!50,dashed] (axis cs:0.602,0) -- (axis cs:0.602,1);
\legend{bank quant,market maker,platform analyst,technology engineer}
\end{axis}
\end{tikzpicture}

\omcaption{Each offer's value as the weight on pay risk varies, the other weights rescaled in proportion. The dashed lines
mark the chosen weight (0.15) and the weight at which the bank's offer overtakes the market maker's (0.602). Data:
\texttt{in\_choose.results} and \texttt{firm.careerdec.flip\_shift}.}\label{fig:in:choosing:sweep}
\end{omfigure}

\section{Revisiting the decision}

A decision about a job is revisited, and the revisiting is where most mistakes happen. Book~16, chapter~26, describes
the tools: a decision journal written at the time, which records the attributes, weights and expectations, and so guards
against outcome bias (judging the choice by how the year went rather than by what was known); and a pre-mortem, which
asks before the move how it could fail. Chapter~28's move model says when leaving is expensive, and chapter~17's \omterm{def:in:quant-researcher:feedback}{feedback
speed} says how long to wait before judging a research job. The model of this chapter is a record of what the chooser
believed; its value is in being rerun when a belief changes.

\section{Tutorial: four offers}\label{tut:in:choosing:tut}

\begin{tutorial}
\textbf{Goal.} Score four offers from the earlier chapters' tools, screen, weight, and measure how robust the choice is.
\textbf{End state:} \cref{fig:in:choosing:ranks,fig:in:choosing:sweep} and the attribute table.
\begin{tutsteps}
\item \textbf{Attributes.} \texttt{firm.payoffer} for pay and its spread; \texttt{firm.locations} for what tax and
housing take; \texttt{firm.workload.oncall\_nights}; judgements for learning; the pay model's closures for security.
\item \textbf{Scale and screen.} \texttt{firm.careerdec.scale} and \texttt{dominated}.
\item \textbf{\omterm{def:in:choosing:smaa}{Rank acceptability}.} \texttt{rank\_acceptability} draws weights on the simplex
(\cref{lst:in:choosing:smaa}).
\omcode{code/firm/careerdec/firm_careerdec.py}{49}{55}{Rank acceptability by sampling weights uniformly on the simplex.}\label{lst:in:choosing:smaa}
\item \textbf{Sensitivity.} \texttt{flip\_shift} searches the smallest change of one weight that changes the best offer
(\cref{lst:in:choosing:flip}).
\omcode{code/firm/careerdec/firm_careerdec.py}{66}{78}{The smallest change of one weight that changes the choice.}\label{lst:in:choosing:flip}
\end{tutsteps}
\end{tutorial}

The results: first-rank acceptability 79.5\% (market maker), 15.9\% (bank quant), 4.6\% (platform analyst), 0\%
(technology engineer); the weight on pay risk must rise by 0.452 to change the choice.

\textbf{What to change next.} Replace the uniform weight distribution with the chooser's ranking of attributes (ordinal
SMAA); add correlation between attributes' uncertainty; run the analysis again after a year with what was learnt.

\section{Build: firm.careerdec}\label{bld:in:choosing:careerdec}

\begin{build}
\bfield{Purpose} Choose among offers with several attributes, and say how robust the choice is to the weights.
\bfield{Interface} \texttt{firm.careerdec}: \texttt{scale(matrix, higher\_better)}; \texttt{swing\_weights(points)};
\texttt{value(scaled, weights)}; \texttt{dominated(matrix, higher\_better)}; \texttt{rank\_acceptability(scaled, n,
rng)}; \texttt{flip\_shift(scaled, weights, attr, step)}; fed by \texttt{firm.payoffer}, \texttt{firm.aftertax},
\texttt{firm.locations}, \texttt{firm.workload} and \texttt{firm.careerpath}.
\bfield{Rules} Scales run from the worst to the best offer; dominance is checked on raw values; weights sum to one; the
flip search changes one weight and rescales the others in proportion.
\bfield{Acceptance tests} \texttt{code/firm/careerdec/tests/}: scaling and reversal; a dominated offer is found; two
mirror-image offers each come first half the time; a flip is found where it must be and not where it cannot be.
\bfield{Stretch} Ordinal weights; interval scores; a value function with interactions between attributes.
\end{build}

\begin{omsources}
\item Keeney, R.~L. and H.~Raiffa (1976), \emph{Decisions with Multiple Objectives: Preferences and Value Tradeoffs},
Wiley (Cambridge University Press, 1993).
\item von Winterfeldt, D. and W.~Edwards (1986), \emph{Decision Analysis and Behavioral Research}, Cambridge University
Press.
\item Lahdelma, R., J.~Hokkanen and P.~Salminen (1998), SMAA---stochastic multiobjective acceptability analysis,
\emph{European Journal of Operational Research} 106(1), 137--143.
\item Chapters~13, 15, 17, 22 and 26--28 of this book; Book~4, chapter~9; Book~16, chapter~26.
\end{omsources}

\section{Exercises}

\begin{exercise}[$\star$]\label{exo:in:choosing:1}
Which offer is dominated, and by which?
\end{exercise}

\begin{exercise}[$\star$]\label{exo:in:choosing:2}
What are the swing weights, as shares, of the six attributes?
\end{exercise}

\begin{exercise}[$\star$]\label{exo:in:choosing:3}
Which offer has the highest certainty equivalent at relative risk aversion 3, and which the highest expected pay?
\end{exercise}

\begin{exercise}[$\star\star$]\label{exo:in:choosing:4}
Why does the platform offer come first in only 4.6\% of weight draws although it ties for the best learning score?
\end{exercise}

\begin{exercise}[$\star\star$]\label{exo:in:choosing:5}
What is the weakness of drawing weights uniformly on the simplex?
\end{exercise}

\begin{exercise}[$\star\star$]\label{exo:in:choosing:6}
A friend weights learning at 60 swing points and the rest as in the chapter. Which offer does the model choose?
\end{exercise}

\begin{exercise}[$\star\star\star$]\label{exo:in:choosing:7}
\emph{Coding.} Drop the dominated offer and recompute the first-rank acceptability of the other three.
\end{exercise}

\begin{exercise}[$\star\star\star$]\label{exo:in:choosing:8}
\emph{Find the flaw.} ``The market maker's offer wins in 79.5\% of weightings, so it is the right choice for 79.5\% of
people.''
\end{exercise}

\section{Problem: Four Offers}

\begin{problem}[{Weekend problem --- four offers}]\label{pb:in:choosing:1}
A graduate holds the four offers of the chapter and asks which to accept.

\medskip\noindent\textbf{Part I --- The attributes.}
\begin{enumerate}
\item Which attributes differ across the offers, and which chapters measure each?
\item Define the \omterm{def:in:choosing:mavm}{multi-attribute value model} and \omterm{def:in:choosing:mavm}{swing weighting}.
\item State the four certainty equivalents at relative risk aversion 3.
\item Why keep expected pay and pay risk apart?
\item Which attributes are judgements?
\end{enumerate}

\medskip\noindent\textbf{Part II --- The model.}
\begin{enumerate}[resume]
\item Give the attribute table.
\item Define a \omterm{def:in:choosing:smaa}{dominated alternative} and name the dominated offer.
\item State the swing points and the offers' values.
\item Define \omterm{def:in:choosing:smaa}{rank acceptability} and state the method.
\item Give the first-rank acceptabilities.
\end{enumerate}

\medskip\noindent\textbf{Part III --- Sensitivity.}
\begin{enumerate}[resume]
\item How far must the weight on pay risk move to change the choice?
\item What does that say about the bank's offer?
\item What would change the learning scores?
\item What would a larger platform payout change?
\item What does the model not know?
\end{enumerate}

\medskip\noindent\textbf{Part IV --- The verdict.}
\begin{enumerate}[resume]
\item State the \emph{named result}: each offer's first-rank acceptability index under uniform weights, and the smallest
change in the weight on pay risk that changes the preferred offer.
\item How should the graduate record the decision?
\item When should she revisit it?
\item What would make her regret it, and how would she tell bad luck from a bad choice?
\item In two sentences, answer her.
\end{enumerate}
\end{problem}

\section{Interview questions}

\begin{interviewq}[$\star$ \iqroles{researcher, trader}]\label{iq:in:choosing:1}
Why do you want this job rather than the others you are considering?
\end{interviewq}

\begin{interviewq}[$\star$ \iqroles{researcher}]\label{iq:in:choosing:2}
How would you compare two offers whose pay differs in structure but not in expected value?
\end{interviewq}

\begin{interviewq}[$\star\star$ \iqroles{researcher, developer}]\label{iq:in:choosing:3}
Write a function that finds which alternatives are dominated. What is its complexity?
\end{interviewq}

\begin{interviewq}[$\star\star$ \iqroles{researcher}]\label{iq:in:choosing:4}
How would you sample weights uniformly on a simplex?
\end{interviewq}

\begin{interviewq}[$\star\star$ \iqroles{trader}]\label{iq:in:choosing:5}
You took a job that went badly in its first year. How do you decide whether it was a bad decision?
\end{interviewq}

\begin{interviewq}[$\star\star\star$ \iqroles{researcher}]\label{iq:in:choosing:6}
Two attributes are strongly correlated across your options. What does that do to an additive value model, and what would
you do about it?
\end{interviewq}
